We thank the reviewer. We replaced the bending range in Section~3.2 with the measured one (single section: maximum $57^\circ$, 95th percentile $36^\circ$; multi-section 95th percentiles of $30^\circ$, $32^\circ$, and $34^\circ$), and we added the coordinates of all targets with their bending angle and length change (Tables~S8 and S9) and the coverage of the evaluation set (SI Section~6.3 and Fig.~S9). The single-section targets and trajectories reach $33.7^\circ$ and $40.9^\circ$, close to the 95th percentile of the training data ($35.6^\circ$); the 20 targets span 9.2\% (single) and 16.4\% (multi) of the convex hull of the sampled positions. We revised the text to state that the reported standard deviations describe the variation across targets visited once, not repeatability, and that trajectory waypoints are serially correlated. The disturbance experiment is now described as the compensation of a single constant 50~g off-center load rather than general robustness. Repeated trials, non-planar and boundary trajectories, and further load conditions are listed as limitations and future work in Section~5.
